\section{Konversi NFA ke DFA: Subset Construction}
\label{sec:ch06-subset}

\subsection{Ide Dasar}

Setiap state DFA merepresentasikan himpunan state NFA yang mungkin
dicapai setelah membaca prefiks tertentu. Jika NFA berada pada
himpunan state \(S\) dan membaca simbol \(a\), state DFA berikutnya
adalah epsilon closure dari semua tujuan transisi \(a\) dari \(S\).
Konstruksi ini disebut \eng{subset construction} atau
\eng{powerset construction}.

\subsection{Algoritma}

\begin{enumerate}[label=\arabic*.,leftmargin=*]
  \item State awal DFA adalah \(\varepsilon\text{-closure}(\{q_0\})\).
  \item Untuk setiap state DFA \(S\) yang belum diproses dan setiap
        simbol \(a \in \Sigma\):
        \begin{enumerate}[label=(\alph*),leftmargin=*]
          \item hitung \(T = \varepsilon\text{-closure}(
                \bigcup_{q \in S} \delta(q,a))\);
          \item jika \(T\) belum ada, tambahkan sebagai state DFA;
          \item tetapkan transisi \(\delta_D(S,a) = T\).
        \end{enumerate}
  \item State DFA menerima jika \(S \cap F \neq \varnothing\).
\end{enumerate}

Jumlah state DFA paling banyak \(2^{|Q|}\), sehingga konversi dapat
meledak secara eksponensial. Dalam praktik, banyak himpunan state
tidak dapat dicapai, sehingga DFA hasil konversi biasanya jauh lebih
kecil.

\subsection{Contoh}

Untuk NFA pada \figref{fig:ch06-nfa-abb}, state awal adalah
\(A = \{0\}\). Transisinya:
\[
\begin{aligned}
A \xrightarrow{a} \{0,1\} &= B, &
A \xrightarrow{b} \{0\} &= A,\\
B \xrightarrow{a} \{0,1\} &= B, &
B \xrightarrow{b} \{0,2\} &= C,\\
C \xrightarrow{a} \{0,1\} &= B, &
C \xrightarrow{b} \{0,3\} &= D,\\
D \xrightarrow{a} \{0,1\} &= B, &
D \xrightarrow{b} \{0\} &= A.
\end{aligned}
\]
State penerima adalah \(D\), karena memuat state NFA 3.

\subsection{Artefak: Tabel Konversi}

Tabel \tabref{tab:ch06-konversi} merangkum hasil subset construction
untuk contoh di atas. Tabel ini adalah artefak Sub-CPMK 2.3: setiap
baris harus diisi dari NFA asal, bukan disalin dari DFA yang sudah
jadi.

\begin{table}[H]
\centering
\caption{Tabel konversi NFA ke DFA untuk contoh pada
\figref{fig:ch06-dfa-abb}}
\label{tab:ch06-konversi}
\begin{tabular}{@{}c l c c c@{}}
\toprule
\textbf{State DFA} & \textbf{Himpunan state NFA} &
\textbf{\(a\)} & \textbf{\(b\)} & \textbf{Penerima} \\
\midrule
\(A\) & \(\{0\}\) & \(B\) & \(A\) & tidak \\
\(B\) & \(\{0,1\}\) & \(B\) & \(C\) & tidak \\
\(C\) & \(\{0,2\}\) & \(B\) & \(D\) & tidak \\
\(D\) & \(\{0,3\}\) & \(B\) & \(A\) & ya \\
\bottomrule
\end{tabular}
\end{table}

Saat mengisi tabel, hitung \(\varepsilon\text{-closure}\) dari
\(\text{move}(S,a)\) untuk setiap himpunan state \(S\). Jika satu
himpunan belum pernah muncul, tambahkan sebagai state DFA baru dan
beri nama yang konsisten.

\begin{figure}[H]
\centering
\begin{tikzpicture}[node distance=22mm]
  \node[state, initial] (A) {\(A\)};
  \node[state, right=of A] (B) {\(B\)};
  \node[state, right=of B] (C) {\(C\)};
  \node[accepting, right=of C] (D) {\(D\)};
  \draw[arrow] (A) edge[loop above] node[font=\small] {\(b\)} (A);
  \draw[arrow] (B) edge[loop above] node[font=\small] {\(a\)} (B);
  \draw[arrow] (A) -- node[above, font=\small] {\(a\)} (B);
  \draw[arrow] (B) -- node[above, font=\small] {\(b\)} (C);
  \draw[arrow] (C) -- node[above, font=\small] {\(b\)} (D);
  \draw[arrow] (C) to[bend left=25] node[below, font=\small] {\(a\)} (B);
  \draw[arrow] (D) to[bend left=40] node[below, font=\small] {\(a\)} (B);
  \draw[arrow] (D) to[bend left=20] node[below, font=\small] {\(b\)} (A);
\end{tikzpicture}
\caption{DFA hasil subset construction dari NFA pada \figref{fig:ch06-nfa-abb}}
\label{fig:ch06-dfa-abb}
\end{figure}

\subsection{Implementasi}

Algoritma subset construction dapat diimplementasikan dengan antrean
state DFA dan kamus transisi. Kunci state DFA adalah
\code{frozenset} dari state NFA. Potongan implementasi Python
ditunjukkan pada \lstref{lst:ch06-to-dfa}.

\includecode[
  language=Python,
  firstline=150,
  lastline=168,
  firstnumber=150,
  caption={Fungsi \code{to\_dfa} pada \file{fa\_tool.py}},
  label={lst:ch06-to-dfa}
]{../../code/tools/fa_tool.py}

Implementasi ini menghasilkan DFA yang mungkin masih memiliki state
tidak terjangkau. Minimalisasi dan penghapusan state tidak
terjangkau dilakukan pada tahap berikutnya.
